\begin{table*}[t]\centering\caption{Main results: $|\hat T_c-T_c|$ (mean $\pm$ std over \rhval{const/n_seeds} seeds) from the single-size read-out at each $L$ and after the FSS collapse, collapse exponent, and held-out far-from-$T_c$ accuracy. Bold best, underline second best (generated).}\label{tab:main}
\begin{tabular}{lccccccc}
\toprule
Method & tc\_error\_l16 $\downarrow$ & tc\_error\_l24 $\downarrow$ & tc\_error\_l32 $\downarrow$ & tc\_error\_fss $\downarrow$ & nu\_error\_fss $\downarrow$ & acc\_far $\uparrow$ & $n$ \\
\midrule
Binder/chi (reference) & 0.109 $\pm$ 0.018 & 0.088 $\pm$ 0.009 & 0.066 $\pm$ 0.007 & 0.035 $\pm$ 0.036 & 0.250 $\pm$ 0.223 & -- & 10 \\
MLP & \underline{0.051 $\pm$ 0.015} & 0.029 $\pm$ 0.010 & \underline{0.020 $\pm$ 0.008} & \textbf{0.021 $\pm$ 0.019} & \underline{0.168 $\pm$ 0.134} & \underline{0.997 $\pm$ 0.001} & 10 \\
LogReg (Z2-fixed) & 0.504 $\pm$ 0.059 & 0.367 $\pm$ 0.054 & 0.287 $\pm$ 0.052 & 0.083 $\pm$ 0.030 & \textbf{0.143 $\pm$ 0.139} & 0.723 $\pm$ 0.011 & 10 \\
Confusion (MLP) & 0.056 $\pm$ 0.052 & \textbf{0.017 $\pm$ 0.017} & \textbf{0.015 $\pm$ 0.016} & 0.331 $\pm$ 0.000 & 1.000 $\pm$ 0.000 & -- & 10 \\
PCA & 0.163 $\pm$ 0.075 & 0.101 $\pm$ 0.047 & 0.056 $\pm$ 0.017 & 0.122 $\pm$ 0.037 & 0.480 $\pm$ 0.265 & -- & 10 \\
LogReg (raw) & 0.409 $\pm$ 0.240 & 0.341 $\pm$ 0.196 & 0.220 $\pm$ 0.218 & 0.155 $\pm$ 0.114 & 0.473 $\pm$ 0.198 & 0.513 $\pm$ 0.011 & 10 \\
\textbf{CNN} & \textbf{0.032 $\pm$ 0.018} & \underline{0.023 $\pm$ 0.009} & 0.023 $\pm$ 0.006 & \underline{0.034 $\pm$ 0.018} & 0.698 $\pm$ 0.299 & \textbf{0.999 $\pm$ 0.001} & 10 \\
\bottomrule
\end{tabular}

\end{table*}

\begin{table}[t]\centering\caption{Signed bias $\hat T_c-T_c$ (mean $\pm$ std over seeds).}\label{tab:bias}
\resizebox{\columnwidth}{!}{\begin{tabular}{lcccc}
\toprule
System & $L{=}16$ & $L{=}24$ & $L{=}32$ & collapse \\
\midrule
CNN & \rhval{main/cnn/ising/tc_bias_l16/mean:3} $\pm$ \rhval{main/cnn/ising/tc_bias_l16/std:3} & \rhval{main/cnn/ising/tc_bias_l24/mean:3} $\pm$ \rhval{main/cnn/ising/tc_bias_l24/std:3} & \rhval{main/cnn/ising/tc_bias_l32/mean:3} $\pm$ \rhval{main/cnn/ising/tc_bias_l32/std:3} & \rhval{main/cnn/ising/tc_bias_fss/mean:3} $\pm$ \rhval{main/cnn/ising/tc_bias_fss/std:3} \\
MLP & \rhval{main/mlp/ising/tc_bias_l16/mean:3} $\pm$ \rhval{main/mlp/ising/tc_bias_l16/std:3} & \rhval{main/mlp/ising/tc_bias_l24/mean:3} $\pm$ \rhval{main/mlp/ising/tc_bias_l24/std:3} & \rhval{main/mlp/ising/tc_bias_l32/mean:3} $\pm$ \rhval{main/mlp/ising/tc_bias_l32/std:3} & \rhval{main/mlp/ising/tc_bias_fss/mean:3} $\pm$ \rhval{main/mlp/ising/tc_bias_fss/std:3} \\
Confusion (MLP) & \rhval{main/confusion-mlp/ising/tc_bias_l16/mean:3} $\pm$ \rhval{main/confusion-mlp/ising/tc_bias_l16/std:3} & \rhval{main/confusion-mlp/ising/tc_bias_l24/mean:3} $\pm$ \rhval{main/confusion-mlp/ising/tc_bias_l24/std:3} & \rhval{main/confusion-mlp/ising/tc_bias_l32/mean:3} $\pm$ \rhval{main/confusion-mlp/ising/tc_bias_l32/std:3} & \rhval{main/confusion-mlp/ising/tc_bias_fss/mean:3} $\pm$ \rhval{main/confusion-mlp/ising/tc_bias_fss/std:3} \\
PCA & \rhval{main/pca/ising/tc_bias_l16/mean:3} $\pm$ \rhval{main/pca/ising/tc_bias_l16/std:3} & \rhval{main/pca/ising/tc_bias_l24/mean:3} $\pm$ \rhval{main/pca/ising/tc_bias_l24/std:3} & \rhval{main/pca/ising/tc_bias_l32/mean:3} $\pm$ \rhval{main/pca/ising/tc_bias_l32/std:3} & \rhval{main/pca/ising/tc_bias_fss/mean:3} $\pm$ \rhval{main/pca/ising/tc_bias_fss/std:3} \\
Binder/$\chi$ (ref.) & \rhval{main/binder-chi-reference/ising/tc_bias_l16/mean:3} $\pm$ \rhval{main/binder-chi-reference/ising/tc_bias_l16/std:3} & \rhval{main/binder-chi-reference/ising/tc_bias_l24/mean:3} $\pm$ \rhval{main/binder-chi-reference/ising/tc_bias_l24/std:3} & \rhval{main/binder-chi-reference/ising/tc_bias_l32/mean:3} $\pm$ \rhval{main/binder-chi-reference/ising/tc_bias_l32/std:3} & \rhval{main/binder-chi-reference/ising/tc_bias_fss/mean:3} $\pm$ \rhval{main/binder-chi-reference/ising/tc_bias_fss/std:3} \\
LogReg (Z2-fixed) & \rhval{main/logreg-z2-fixed/ising/tc_bias_l16/mean:3} $\pm$ \rhval{main/logreg-z2-fixed/ising/tc_bias_l16/std:3} & \rhval{main/logreg-z2-fixed/ising/tc_bias_l24/mean:3} $\pm$ \rhval{main/logreg-z2-fixed/ising/tc_bias_l24/std:3} & \rhval{main/logreg-z2-fixed/ising/tc_bias_l32/mean:3} $\pm$ \rhval{main/logreg-z2-fixed/ising/tc_bias_l32/std:3} & \rhval{main/logreg-z2-fixed/ising/tc_bias_fss/mean:3} $\pm$ \rhval{main/logreg-z2-fixed/ising/tc_bias_fss/std:3} \\
LogReg (raw) & \rhval{main/logreg-raw/ising/tc_bias_l16/mean:3} $\pm$ \rhval{main/logreg-raw/ising/tc_bias_l16/std:3} & \rhval{main/logreg-raw/ising/tc_bias_l24/mean:3} $\pm$ \rhval{main/logreg-raw/ising/tc_bias_l24/std:3} & \rhval{main/logreg-raw/ising/tc_bias_l32/mean:3} $\pm$ \rhval{main/logreg-raw/ising/tc_bias_l32/std:3} & \rhval{main/logreg-raw/ising/tc_bias_fss/mean:3} $\pm$ \rhval{main/logreg-raw/ising/tc_bias_fss/std:3} \\
\bottomrule
\end{tabular}
}
\end{table}

\begin{figure*}[t]\centering\includegraphics[width=\textwidth]{fig_curves.pdf}
\caption{Mean held-out output $P(\text{ordered})$ versus $T$ (band: std over seeds; dotted: exact $T_c$). The raw logistic regression is flat at $1/2$.}\label{fig:curves}\end{figure*}

\subsection{H1: nonlinear classifiers versus PCA at a single size --- supported}
At $L=32$ the CNN error is \rhval{main/cnn/ising/tc_error_l32/mean:3} and the MLP error \rhval{main/mlp/ising/tc_error_l32/mean:3}, versus \rhval{main/pca/ising/tc_error_l32/mean:3} for PCA and \rhval{main/binder-chi-reference/ising/tc_error_l32/mean:3} for the susceptibility peak (Table~\ref{tab:main}). The paired Wilcoxon tests give $p=\rhval{analysis/paired-tests/ising/h1_cnn_vs_pca_l32_p/mean:3}$ (CNN) and $p=\rhval{analysis/paired-tests/ising/h1_mlp_vs_pca_l32_p/mean:3}$ (MLP), with the CNN lower on \rhval{analysis/paired-tests/ising/h1_cnn_vs_pca_l32_wins/mean:0} of \rhval{analysis/paired-tests/ising/n_paired_seeds/mean:0} seeds and the MLP on \rhval{analysis/paired-tests/ising/h1_mlp_vs_pca_l32_wins/mean:0}; the paired $t$-test of \texttt{rh compare} agrees for the CNN ($p=\rhval{cmp/main/pca/ising/tc_error_l32/paired_p:4}$). The advantage is largest at the small sizes (Table~\ref{tab:main}). The result does not extend to confusion: its $L=32$ error (\rhval{main/confusion-mlp/ising/tc_error_l32/mean:3}) is nominally the lowest of all systems, but not distinguishable from the CNN (Wilcoxon $p=\rhval{analysis/paired-tests/ising/conf_vs_cnn_l32_p/mean:2}$) or the MLP ($p=\rhval{analysis/paired-tests/ising/conf_vs_mlp_l32_p/mean:2}$). At $L=24$ it is nominally lowest too (\rhval{main/confusion-mlp/ising/tc_error_l24/mean:3}; $p=\rhval{analysis/paired-tests/ising/conf_vs_cnn_l24_p/mean:2}$ against the CNN and $\rhval{analysis/paired-tests/ising/conf_vs_mlp_l24_p/mean:3}$ against the MLP, uncorrected for the six comparisons made at the three sizes), and its seed spread is the largest of the learned systems at $L=16$ (Table~\ref{tab:main}). Note that all errors for the good systems are below the grid spacing of \rhval{const/temp_step:3}, so they are interpolation-limited, and the $L=32$ error difference between CNN and MLP is not significant ($p=\rhval{analysis/paired-tests/ising/cnn_vs_mlp_l32_p/mean:2}$).

\emph{Finite-size bias.} Table~\ref{tab:bias} shows that the single-size read-outs are biased upward and the bias shrinks with $L$ for the MLP (\rhval{main/mlp/ising/tc_bias_l16/mean:3}, \rhval{main/mlp/ising/tc_bias_l24/mean:3}, \rhval{main/mlp/ising/tc_bias_l32/mean:3} at $L=16,24,32$), PCA and the susceptibility peak, as expected from the shift of pseudo-critical points with $L$~\cite{fisher1972scaling}. The CNN is the exception: its bias is small already at $L=16$ (\rhval{main/cnn/ising/tc_bias_l16/mean:3}) and nearly constant afterwards. We did not test why; the CNN output curves are nearly identical for the three sizes (Fig.~\ref{fig:curves}), which is consistent with a read-out dominated by local features whose statistics change little with $L$, but this is a conjecture.

\emph{The CNN is not the best system.} The platform's verdict on the primary metric (collapse error) ranks the MLP first by mean (\rhval{main/mlp/ising/tc_error_fss/mean:3} versus \rhval{main/cnn/ising/tc_error_fss/mean:3}), but the difference is not significant (Table~\ref{tab:tests}); the CNN is nominally best only at $L=16$, learning by confusion at $L=24$ and $L=32$, and no learned system is clearly distinguishable from the others at $L=32$.

\subsection{H2: collapse versus single-size crossing --- not supported}
The collapse does not reduce the error (Table~\ref{tab:tests}). For the CNN the collapse error (\rhval{main/cnn/ising/tc_error_fss/mean:3}) is larger than at $L=32$ (mean difference \rhval{analysis/paired-tests/ising/h2_cnn_fss_vs_l32_meandiff/mean:3}, $p=\rhval{analysis/paired-tests/ising/h2_cnn_fss_vs_l32_p/mean:3}$, lower on only \rhval{analysis/paired-tests/ising/h2_cnn_fss_vs_l32_wins/mean:0} seeds); for the MLP there is no difference ($p=\rhval{analysis/paired-tests/ising/h2_mlp_fss_vs_l32_p/mean:2}$). The collapse helps where the single-size bias is large and curves differ in position only: the susceptibility reference improves from \rhval{main/binder-chi-reference/ising/tc_error_l32/mean:3} to \rhval{main/binder-chi-reference/ising/tc_error_fss/mean:3} with Binder collapse ($p=\rhval{analysis/paired-tests/ising/binder_fss_vs_l32_p/mean:3}$, which is not below $0.05$), whereas PCA worsens (\rhval{main/pca/ising/tc_error_l32/mean:3} to \rhval{main/pca/ising/tc_error_fss/mean:3}, $p=\rhval{analysis/paired-tests/ising/pca_fss_vs_l32_p/mean:3}$). Statistically, with only ten seeds the smallest attainable two-sided Wilcoxon $p$ is $0.002$, so ``not supported'' means no detectable improvement, not proof of none.

\begin{table}[t]\centering\caption{Paired Wilcoxon tests over seeds (second column: first minus second; ``seeds first lower'' counts out of ten).}\label{tab:tests}
\resizebox{\columnwidth}{!}{\begin{tabular}{llccc}
\toprule
Hyp. & comparison (first minus second) & mean diff. & Wilcoxon $p$ & seeds first lower \\
\midrule
H1 & CNN vs PCA, $L{=}32$ & \rhval{analysis/paired-tests/ising/h1_cnn_vs_pca_l32_meandiff/mean:3} & \rhval{analysis/paired-tests/ising/h1_cnn_vs_pca_l32_p/mean:3} & \rhval{analysis/paired-tests/ising/h1_cnn_vs_pca_l32_wins/mean:0} \\
H1 & MLP vs PCA, $L{=}32$ & \rhval{analysis/paired-tests/ising/h1_mlp_vs_pca_l32_meandiff/mean:3} & \rhval{analysis/paired-tests/ising/h1_mlp_vs_pca_l32_p/mean:3} & \rhval{analysis/paired-tests/ising/h1_mlp_vs_pca_l32_wins/mean:0} \\
H2 & CNN: collapse vs $L{=}32$ & \rhval{analysis/paired-tests/ising/h2_cnn_fss_vs_l32_meandiff/mean:3} & \rhval{analysis/paired-tests/ising/h2_cnn_fss_vs_l32_p/mean:3} & \rhval{analysis/paired-tests/ising/h2_cnn_fss_vs_l32_wins/mean:0} \\
H2 & MLP: collapse vs $L{=}32$ & \rhval{analysis/paired-tests/ising/h2_mlp_fss_vs_l32_meandiff/mean:3} & \rhval{analysis/paired-tests/ising/h2_mlp_fss_vs_l32_p/mean:3} & \rhval{analysis/paired-tests/ising/h2_mlp_fss_vs_l32_wins/mean:0} \\
H4 & PCA vs CNN (collapse) & \rhval{analysis/paired-tests/ising/h4_pca_vs_cnn_fss_meandiff/mean:3} & \rhval{analysis/paired-tests/ising/h4_pca_vs_cnn_fss_p/mean:3} & \rhval{analysis/paired-tests/ising/h4_pca_vs_cnn_fss_wins/mean:0} \\
H4 & Confusion vs CNN (collapse) & \rhval{analysis/paired-tests/ising/h4_conf_vs_cnn_fss_meandiff/mean:3} & \rhval{analysis/paired-tests/ising/h4_conf_vs_cnn_fss_p/mean:3} & \rhval{analysis/paired-tests/ising/h4_conf_vs_cnn_fss_wins/mean:0} \\
-- & Confusion vs CNN, $L{=}32$ & \rhval{analysis/paired-tests/ising/conf_vs_cnn_l32_meandiff/mean:3} & \rhval{analysis/paired-tests/ising/conf_vs_cnn_l32_p/mean:3} & \rhval{analysis/paired-tests/ising/conf_vs_cnn_l32_wins/mean:0} \\
-- & Confusion vs MLP, $L{=}32$ & \rhval{analysis/paired-tests/ising/conf_vs_mlp_l32_meandiff/mean:3} & \rhval{analysis/paired-tests/ising/conf_vs_mlp_l32_p/mean:3} & \rhval{analysis/paired-tests/ising/conf_vs_mlp_l32_wins/mean:0} \\
-- & CNN: $\nu{=}1$ vs free (collapse) & \rhval{analysis/paired-tests/ising/nu1_cnn_meandiff/mean:3} & \rhval{analysis/paired-tests/ising/nu1_cnn_p/mean:3} & \rhval{analysis/paired-tests/ising/nu1_cnn_wins/mean:0} \\
-- & MLP: $\nu{=}1$ vs free (collapse) & \rhval{analysis/paired-tests/ising/nu1_mlp_meandiff/mean:3} & \rhval{analysis/paired-tests/ising/nu1_mlp_p/mean:3} & \rhval{analysis/paired-tests/ising/nu1_mlp_wins/mean:0} \\
-- & MLP vs CNN (collapse) & \rhval{analysis/paired-tests/ising/h4_mlp_vs_cnn_fss_meandiff/mean:3} & \rhval{analysis/paired-tests/ising/h4_mlp_vs_cnn_fss_p/mean:3} & \rhval{analysis/paired-tests/ising/h4_mlp_vs_cnn_fss_wins/mean:0} \\
-- & Binder/$\chi$ vs CNN (collapse) & \rhval{analysis/paired-tests/ising/h4_binder_vs_cnn_fss_meandiff/mean:3} & \rhval{analysis/paired-tests/ising/h4_binder_vs_cnn_fss_p/mean:3} & \rhval{analysis/paired-tests/ising/h4_binder_vs_cnn_fss_wins/mean:0} \\
\bottomrule
\end{tabular}
}
\end{table}

\subsection{H3: logistic regression and the spin-flip symmetry --- supported}
The raw logistic regression gives held-out accuracy \rhval{main/logreg-raw/ising/acc_far/mean:2} inside the training windows (chance level for the class-balanced task), a flat output (Fig.~\ref{fig:curves}), and a crossing that is noise: its error is \rhval{main/logreg-raw/ising/tc_error_l32/mean:2} at $L=32$ with std \rhval{main/logreg-raw/ising/tc_error_l32/std:2}, and a bias with either sign. A linear function of $s$ has zero average over a symmetric ensemble, so it cannot separate the phases. Multiplying each configuration by the sign of its magnetisation, the only change in that system, raises the held-out accuracy to \rhval{main/logreg-z2-fixed/ising/acc_far/mean:2} ($>0.6$, the registered threshold), which isolates the symmetry as the cause. The recovery is partial: the classifier now responds to $|m|$ but its crossing is badly biased (\rhval{main/logreg-z2-fixed/ising/tc_bias_l32/mean:2} at $L=32$, \rhval{main/logreg-z2-fixed/ising/tc_bias_l16/mean:2} at $L=16$) possibly because a linear function of the aligned spins separates the phases poorly near the transition; we did not test this explanation.

\subsection{H4: label-free methods versus the CNN --- supported for PCA; for confusion supported only through a failed collapse}
On the collapse error PCA (\rhval{main/pca/ising/tc_error_fss/mean:3}) and confusion (\rhval{main/confusion-mlp/ising/tc_error_fss/mean:3}) are both worse than the CNN (\rhval{main/cnn/ising/tc_error_fss/mean:3}), $p=\rhval{analysis/paired-tests/ising/h4_pca_vs_cnn_fss_p/mean:3}$ and $p=\rhval{analysis/paired-tests/ising/h4_conf_vs_cnn_fss_p/mean:3}$, with the CNN lower on every seed (Table~\ref{tab:tests}). The confusion value is a \emph{failure case}: for every seed the collapse returns the same grid-boundary value, $T_c=2.6$ with $\nu=2$ (std \rhval{main/confusion-mlp/ising/tc_error_fss/std:3}). This is not just the feasible range of the window: the trial boundaries run from $1.7$ to $3.3$, which with a window of $0.6$ excludes $T_c$ itself, but the failure persists with the narrower windows $0.45$ and $0.3$ that include it (Table~\ref{tab:collapse}). The W-shaped accuracy curves are not functions of $(T-T_c)L^{1/\nu}$ in the sense the collapse cost needs; their \emph{single-size} read-out (central peak) is fine (Table~\ref{tab:main}). MLP and the Binder/susceptibility reference are not distinguishable from the CNN on the collapse error ($p=\rhval{analysis/paired-tests/ising/h4_mlp_vs_cnn_fss_p/mean:2}$ and $\rhval{analysis/paired-tests/ising/h4_binder_vs_cnn_fss_p/mean:2}$). (For confusion this verdict says that our collapse recipe fails, not that its single-size read-out is worse: that read-out is nominally better than the CNN's at $L=24$ and $L=32$, H1 paragraph.) PCA works at a single size (\rhval{main/pca/ising/tc_error_l32/mean:3} at $L=32$, but \rhval{main/pca/ising/tc_error_l16/mean:3} at $L=16$ with a large seed spread, std \rhval{main/pca/ising/tc_error_l16/std:3}).

\begin{figure}[t]\centering\includegraphics[width=\columnwidth]{fig_confusion_pca.pdf}
\caption{Left: learning-by-confusion accuracy $W(T')$ (mean over seeds). Right: PCA order-parameter proxy. Dotted: exact $T_c$.}\label{fig:conf}\end{figure}

\subsection{H5: collapse exponent --- mixed}
The MLP collapse exponent is $\hat\nu=\rhval{main/mlp/ising/nu_fss/mean:2}$ ($|\hat\nu-1|=\rhval{main/mlp/ising/nu_error_fss/mean:2}$, below the registered $0.25$), but for the CNN $\hat\nu=\rhval{main/cnn/ising/nu_fss/mean:2}$ with $|\hat\nu-1|=\rhval{main/cnn/ising/nu_error_fss/mean:2}$ and the search reaches the upper grid edge (maximum over seeds \rhval{main/cnn/ising/nu_fss/max:2}); the verdict is supported for the MLP and refuted for the CNN. Raw logistic regression has $|\hat\nu-1|=\rhval{main/logreg-raw/ising/nu_error_fss/mean:2}$ with a collapse cost far larger (\rhval{main/logreg-raw/ising/collapse_cost/mean:2} versus \rhval{main/mlp/ising/collapse_cost/mean:4}), so its exponent should be read as noise. Fig.~\ref{fig:collapse} shows seed-0 collapses; the CNN curves are nearly size-independent, which makes $\nu$ poorly determined. Fixing $\nu=1$ changes the CNN error to \rhval{abl_nu_fixed/cnn/ising/tc_error_fss/mean:3} (Table~\ref{tab:collapse}).

\begin{figure*}[t]\centering\includegraphics[width=\textwidth]{fig_collapse.pdf}
\caption{Collapse of the output curves with the $(T_c,\nu)$ fitted for seed 0. Binder cumulant panel: reference system.}\label{fig:collapse}\end{figure*}

\begin{figure}[t]\centering\includegraphics[width=\columnwidth]{bars_main_tc_error_fss.pdf}
\caption{Collapse error $|\hat T_c^{\rm fss}-T_c|$ by system (bars: mean over seeds).}\label{fig:bars}\end{figure}
